\documentclass[11pt]{article}

\usepackage[T1]{fontenc}
\usepackage{amsmath,amssymb,amsthm}
\usepackage[margin=1in]{geometry}
\usepackage[hidelinks]{hyperref}
\setlength{\emergencystretch}{3em}

\newtheorem{theorem}{Theorem}
\newtheorem{lemma}[theorem]{Lemma}
\newtheorem{corollary}[theorem]{Corollary}
\newtheorem{conjecture}[theorem]{Conjecture}
\theoremstyle{remark}
\newtheorem{remark}[theorem]{Remark}

\newcommand{\R}{\mathbb R}

\title{The Temporal Weak-\(L^3\) Five-Power Barrier\\
and Exclusion of Retained Fast Type-II Carriers}
\author{John Robert Lawson and Codex}
\date{Research round ending 24 July 2026}

\begin{document}
\maketitle

\begin{abstract}
This round closes the retained partial/tight cell of the
energy-efficient exact \(q=4\) Type-II ledger.  The record clock puts the
weak-\(L^3\) norm in \(L^s\) in time for every \(s<11/2\).  Interpolation
with the energy-class \(L^2_tL^6_x\) bound produces a strong
\(L^q_tL^p_x\) pair strictly inside a peer-reviewed
Leslie--Shvydkoy positive-energy-dimension region.  A fixed-energy
shrinking carrier would create an atom in the terminal energy measure,
whereas that theorem forbids atoms.  With the explicit choice
\[
 s=\frac{16}{3},\qquad
 p=\frac{96}{31},\qquad
 q=\frac{256}{53},
\]
the published local-energy exponents are \(1/29\) in space and \(72/29\)
in terminal time.  The exact \(q=4\) carrier floor contradicts the
resulting decay directly.  The diffuse and
divergent-normalised-energy branches remain open, and the Clay problem is
unsolved.
\end{abstract}

\section{Epistemic scope}

The imported energy-measure estimate is a peer-reviewed theorem of
Leslie and Shvydkoy~\cite{LS18}.  The record-clock conversion, Lorentz
interpolation, tightness lemma, explicit exponent bridge, and application
to the repository's exact \(q=4\) ledger are deductions in this round and
await external mathematical review.  No claim of literature novelty is
made.

The conclusion is a conditional exclusion theorem: no smooth
finite-energy trajectory can satisfy all hypotheses defining the retained
energy-efficient exact \(q=4\) cell.  It is neither a global regularity
theorem nor a breakdown construction.

\section{First-blow-up setting}

Let \(u\) be a smooth divergence-free solution of
\[
 \partial_tu+u\cdot\nabla u=-\nabla p+\nu\Delta u,
 \qquad
 \nabla\cdot u=0
\]
on \(\R^3\times[0,T^*)\), where \(T^*<\infty\) is the first possible
singular time.  Assume
\[
 \sup_{t<T^*}\frac12\|u(t)\|_2^2\le E_0,
 \qquad
 \nu\int_0^{T^*}\|\nabla u(t)\|_2^2\,dt\le E_0.
\]
A Leray--Hopf continuation through \(T^*\), obtained by restarting at a
regular earlier time and using weak--strong uniqueness before \(T^*\),
places the solution in the setting of~\cite{LS18}.

Define
\[
 M(t):=\|u(t)\|_{L^{3,\infty}(\R^3)}.
\]
Let \(t_j\uparrow T^*\) be first record times and put
\[
 m_j:=M(t_j).
\]
Thus
\begin{equation}
 \label{eq:first-record}
 M(t)\le m_j\qquad(0\le t\le t_j).
\end{equation}

The exact \(q=4\) ledger has
\begin{equation}
 \label{eq:q4}
 m_j\asymp2^{2j},
 \qquad
 t_{j+1}-t_j\asymp\frac{2^{-11j}}j.
\end{equation}
In its retained energy-efficient geometry cell there are fixed
\(A<\infty\), \(\gamma>0\), centres \(x_j\), and radii
\begin{equation}
 \label{eq:radii}
 R_j\asymp m_j^{-2}\asymp2^{-4j}
\end{equation}
such that
\begin{equation}
 \label{eq:floor}
 \int_{B_{AR_j}(x_j)}|u(x,t_j)|^2\,dx\ge\gamma.
\end{equation}

\section{The imported energy-measure theorem}

The terminal energy measure is
\[
 \mathcal E
 :=
 \underset{t\uparrow T^*}{\operatorname{w^*\!-\!lim}}
 |u(t)|^2\,dx.
\]
The local energy identity makes this limit unique on bounded spatial
domains.

\begin{theorem}[Leslie--Shvydkoy, specialised to three-dimensional NSE]
\label{thm:LS}
Suppose
\[
 u\in L^q(0,T^*;L^p(\R^3)),
 \qquad
 3\le p<\infty,
 \qquad
 \frac6p+\frac5q\le3.
\]
Set
\[
 \alpha
 :=
 \frac{q}{q-1}\left(1+\frac3p\right),
 \qquad
 \beta
 :=
 \frac{q}{q-1}
 \left(3-\frac6p-\frac5q\right).
\]
For every compact \(K\subset\R^3\), there are \(r_K>0\) and \(C_K<\infty\)
such that
\begin{equation}
 \label{eq:LS-local}
 \sup_{\substack{T^*-r^\alpha<t<T^*\\x_0\in K}}
 \int_{B_r(x_0)}|u(x,t)|^2\,dx
 \le C_Kr^\beta
 \qquad(0<r<r_K).
\end{equation}
In particular, if \(6/p+5/q<3\), then the terminal energy measure has no
atoms.
\end{theorem}

\begin{remark}
The paper states the Euler local-dimension proposition and then proves
that it remains valid for three-dimensional Navier--Stokes solutions at a
first blow-up time.  Strong \(L^q_tL^p_x\) is more than the weak-in-time
hypothesis allowed for finite \(p\).

The accepted source prints \(\beta>2\) at one viscous comparison where its
displayed Prodi--Serrin equivalence and the comparison of \(r^{-2}\) with
the time-cutoff power require \(\alpha>2\).  Its following proposition
states the Navier--Stokes extension.  For the explicit pair used below,
\[
 \frac3p+\frac2q=\frac{177}{128}>1,
 \qquad
 \alpha=\frac{72}{29}>2,
\]
so that comparison is checked directly.
\end{remark}

\section{The record clock gives more than five time powers}

\begin{lemma}[Clock-to-integrability conversion]
\label{lem:clock}
Under \eqref{eq:first-record} and \eqref{eq:q4},
\[
 M\in L^s(0,T^*)
 \qquad\text{for every }s<\frac{11}{2}.
\]
\end{lemma}

\begin{proof}
The solution is smooth away from \(T^*\), so only a terminal tail matters.
On \([t_j,t_{j+1}]\), the first-record property at \(t_{j+1}\) gives
\[
 M(t)\le m_{j+1}.
\]
Consequently,
\[
\begin{aligned}
 \int_{t_J}^{T^*}M(t)^s\,dt
 &\le
 \sum_{j\ge J}
 m_{j+1}^s(t_{j+1}-t_j)\\
 &\le
 C_s\sum_{j\ge J}
 \frac{2^{(2s-11)j}}j.
\end{aligned}
\]
This converges for \(s<11/2\).  At \(s=11/2\), this estimate becomes
harmonic and makes no endpoint claim.
\end{proof}

\section{The temporal five-power interpolation barrier}

\begin{lemma}[Lorentz--Sobolev interpolation]
\label{lem:interpolation}
Assume \(M\in L^s(0,T^*)\) for some \(s>5\).  For \(0<\theta<1\), define
\begin{equation}
 \label{eq:pq-def}
 \frac1p=\frac{1-\theta}{3}+\frac{\theta}{6},
 \qquad
 \frac1q=\frac{1-\theta}{s}+\frac{\theta}{2}.
\end{equation}
Then \(u\in L^q_tL^p_x\), and
\begin{equation}
 \label{eq:gap}
 3-\frac6p-\frac5q
 =
 1-\frac5s
 -
 \theta\left(\frac32-\frac5s\right).
\end{equation}
For all sufficiently small \(\theta>0\), the right-hand side is positive.
\end{lemma}

\begin{proof}
Real interpolation gives
\[
 \|u(t)\|_p
 \le
 C
 \|u(t)\|_{L^{3,\infty}}^{1-\theta}
 \|u(t)\|_6^\theta.
\]
Sobolev and the energy identity give
\[
 \|u(t)\|_6\le C\|\nabla u(t)\|_2,
 \qquad
 \nabla u\in L^2_tL^2_x.
\]
H\"older's inequality in time, using \eqref{eq:pq-def}, yields
\[
 \|u\|_{L^q_tL^p_x}
 \le
 C
 \|M\|_{L^s_t}^{1-\theta}
 \|\nabla u\|_{L^2_tL^2_x}^{\theta}.
\]
Finally \(6/p=2-\theta\); substitution gives \eqref{eq:gap}.  Since
\(s>5\), its value at \(\theta=0\) is strictly positive.
\end{proof}

\begin{theorem}[Temporal five-power no-carrier barrier]
\label{thm:five}
If
\[
 M\in L^s(0,T^*)
 \qquad\text{for some }s>5,
\]
there do not exist \(t_j\uparrow T^*\), \(r_j\downarrow0\), \(x_j\in\R^3\),
and \(\gamma>0\) such that
\[
 \int_{B_{r_j}(x_j)}|u(x,t_j)|^2\,dx\ge\gamma
\]
for every \(j\).
\end{theorem}

\begin{proof}
Choose \(\theta>0\) sufficiently small in Lemma~\ref{lem:interpolation}.
Then Theorem~\ref{thm:LS} applies with \(\beta>0\), and therefore
\begin{equation}
 \label{eq:no-atoms}
 \mathcal E(\{x\})=0
 \qquad(x\in\R^3).
\end{equation}

Lemma~\ref{lem:tightness} below makes the family
\(\{|u(t)|^2\,dx:t<T^*\}\) uniformly tight.  Hence the centres of any
fixed-energy shrinking balls are bounded.  Pass to a subsequence with
\(x_j\to x_*\).  For every \(\rho>0\),
\[
 B_{r_j}(x_j)\subset\overline B_\rho(x_*)
\]
eventually.  The closed-set Portmanteau inequality gives
\[
 \mathcal E(\overline B_\rho(x_*))
 \ge
 \limsup_{j\to\infty}
 \int_{\overline B_\rho(x_*)}|u(x,t_j)|^2\,dx
 \ge\gamma.
\]
Continuity from above as \(\rho\downarrow0\) implies
\[
 \mathcal E(\{x_*\})\ge\gamma,
\]
contradicting \eqref{eq:no-atoms}.
\end{proof}

\begin{remark}[Sharpness of the method]
At \(s=5\), the first positive term on the right of \eqref{eq:gap}
vanishes and every \(\theta>0\) makes the gap negative.  Thus five is the
exact temporal threshold of this interpolation-to-energy-dimension route,
not an arbitrary choice.
\end{remark}

\section{Uniform spatial tightness}

\begin{lemma}[Uniform energy tightness]
\label{lem:tightness}
For the finite-energy trajectory above,
\[
 \lim_{R\to\infty}
 \sup_{0\le t<T^*}
 \int_{|x|>R}|u(x,t)|^2\,dx
 =0.
\]
\end{lemma}

\begin{proof}
Choose
\[
 p=\mathcal R_a\mathcal R_b(u_au_b).
\]
Interpolation between \(L^2\) and \(L^6\), followed by H\"older in time,
gives
\[
 \int_0^{T^*}\|u(t)\|_3^3\,dt<\infty.
\]
The Riesz-transform estimate
\[
 \|p(t)\|_{3/2}\le C\|u(t)\|_3^2
\]
therefore gives
\begin{equation}
 \label{eq:flux-action}
 \mathcal F_*
 :=
 \int_0^{T^*}\!\!\int_{\R^3}
 \left(|u|^3+|p||u|\right)\,dx\,dt
 <\infty.
\end{equation}

Let \(0\le\psi_R\le1\) vanish on \(B_R\), equal one outside \(B_{2R}\),
and obey
\[
 \|\nabla\psi_R\|_\infty\le\frac CR,
 \qquad
 \|\Delta\psi_R\|_\infty\le\frac C{R^2}.
\]
Approximate \(\psi_R\) by compactly supported cutoffs and pass to the limit
using finite energy and \eqref{eq:flux-action}.  The local energy equality,
after dropping its nonnegative dissipation term, gives
\[
\begin{aligned}
 \sup_{t<T^*}
 \int_{|x|>2R}|u(x,t)|^2\,dx
 \le{}&
 C\int_{|x|>R}|u_0(x)|^2\,dx\\
 &+\frac{C\mathcal F_*}{R}
 +\frac{C\nu T^*}{R^2}
 \sup_{t<T^*}\|u(t)\|_2^2.
\end{aligned}
\]
Every term tends to zero as \(R\to\infty\).
\end{proof}

\section{Exact q-four exclusion}

Lemma~\ref{lem:clock} allows the explicit choice
\[
 s=\frac{16}{3},
 \qquad
 \theta=\frac1{16}.
\]
Equations \eqref{eq:pq-def} then give
\[
 \frac1p=\frac{31}{96},
 \qquad
 \frac1q=\frac{53}{256},
\]
so
\begin{equation}
 \label{eq:explicit-pq}
 p=\frac{96}{31},
 \qquad
 q=\frac{256}{53}.
\end{equation}
The exact published admissibility gap is
\[
 3-\frac6p-\frac5q
 =
 \frac7{256}>0.
\]
Consequently,
\begin{equation}
 \label{eq:explicit-exponents}
 \beta
 =
 \frac{q}{q-1}
 \left(3-\frac6p-\frac5q\right)
 =
 \frac1{29},
 \qquad
 \alpha
 =
 \frac{q}{q-1}
 \left(1+\frac3p\right)
 =
 \frac{72}{29}.
\end{equation}

\begin{corollary}[Retained exact \(q=4\) cell is empty]
\label{cor:q4}
No smooth finite-energy trajectory can satisfy simultaneously
\eqref{eq:q4}, \eqref{eq:radii}, and \eqref{eq:floor}.
\end{corollary}

\begin{proof}
This follows immediately from Theorem~\ref{thm:five}.  There is also a
direct quantitative contradiction.

The geometric gap sum in \eqref{eq:q4} gives
\[
 T^*-t_j\le C\frac{2^{-11j}}j.
\]
Using \eqref{eq:radii} and \eqref{eq:explicit-exponents},
\[
 \frac{T^*-t_j}{R_j^{72/29}}
 \le
 \frac Cj
 2^{-\left(11-\frac{288}{29}\right)j}
 =
 \frac Cj2^{-31j/29}
 \longrightarrow0.
\]
Lemma~\ref{lem:tightness} places all retained centres in one compact set.
Theorem~\ref{thm:LS}, applied with \(r=AR_j\), therefore gives
\[
 \int_{B_{AR_j}(x_j)}|u(x,t_j)|^2\,dx
 \le
 C(AR_j)^{1/29}
 \longrightarrow0,
\]
contradicting \eqref{eq:floor}.
\end{proof}

\begin{remark}
The preceding terminal-infrared and multirecord-import results remain valid
conditional theorems.  Corollary~\ref{cor:q4} shows that all their retained
exact \(q=4\) premises cannot occur together on an actual smooth
trajectory.  No packet-genealogy argument is needed inside an empty cell.
\end{remark}

\section{Optimised dimension ledger}

The explicit pair \eqref{eq:explicit-pq} is enough for exclusion.  The full
clock gives a sharper quantitative terminal statement.

Fix \(5<s<11/2\) and let \(\theta\downarrow0\) in
Lemma~\ref{lem:interpolation}.  The lower local-dimension exponent tends to
\[
 \frac{s-5}{s-1}.
\]
For sufficiently small \(\theta\), the pair \((p,q)\) lies in Region II of
the refined Leslie--Shvydkoy concentration theorem.  Its lower bound tends
to
\[
 3-\frac{2/s}{1-2/3-1/s}
 =
 \frac{3(s-5)}{s-3}.
\]
Letting \(s\uparrow11/2\) gives
\begin{equation}
 \label{eq:dimensions}
 \boxed{
 d(x,\mathcal E)\ge\frac19
 \quad\text{for every }x,
 \qquad
 D(\mathcal E)\ge\frac35.
 }
\end{equation}
These bounds are not needed for Corollary~\ref{cor:q4}; they record how far
the exact clock lies inside the no-atom regime.

\section{Findings and remaining obligations}

\subsection*{Robust findings, subject to review}

\begin{enumerate}
\item Temporal \(L^s_tL^{3,\infty}_x\) control with any \(s>5\), combined
with the energy class, excludes every fixed-energy shrinking carrier at a
first blow-up time.
\item The number five is the exact endpoint of this interpolation route.
\item The exact \(q=4\) ledger supplies every \(s<11/2\), so it lies
strictly beyond the barrier.
\item The explicit pair \((p,q)=(96/31,256/53)\) gives the positive spatial
power \(1/29\) and terminal window power \(72/29\).
\item The retained partial/tight energy-efficient exact \(q=4\) cell is
empty.
\item The optimised terminal energy measure has lower local dimension at
least \(1/9\) and concentration dimension at least \(3/5\).
\end{enumerate}

\subsection*{Things still to prove}

\begin{enumerate}
\item Exclude or rigidify the diffuse energy-efficient selected layer.
\item Derive an interface, enstrophy, or pressure-transport charge from
fragmentation into many subcarrier components.
\item Control the divergent-normalised-energy branch.
\item Treat other Type-II clocks that fail every temporal
\(L^sL^{3,\infty}\) condition with \(s>5\), especially the endpoint \(s=5\).
\item Prove regularity or construct admissible breakdown for all Clay data.
\end{enumerate}

\begin{conjecture}[Diffuse-layer fragmentation cost]
Suppose an energy-efficient amplitude layer has fixed energy, total volume
comparable to \(R_j^3\), and vanishing energy fraction in every
radius-\(R_j\) ball.  Then either a recentered sublayer has a fixed local
energy floor, or the layer pays a scale-zero nonsummable perimeter,
enstrophy, or pressure-transport charge along an infinite first-record
sequence.
\end{conjecture}

No part of this conjecture is proved here.  This round establishes no Clay
alternative A--D.

\begin{thebibliography}{9}

\bibitem{LS18}
T.~M. Leslie and R. Shvydkoy,
\emph{The Energy Measure for the Euler and Navier--Stokes Equations},
Archive for Rational Mechanics and Analysis \textbf{230} (2018),
459--492.
\href{https://doi.org/10.1007/s00205-018-1250-4}
{doi:10.1007/s00205-018-1250-4}.
Accepted source: arXiv:1705.04420v4.

\end{thebibliography}

\end{document}
